\section{Analisi sperimentale e lezioni apprese}
\label{sec:lessons}

Questa sezione documenta il \emph{processo} sperimentale, non solo i risultati
finali: cosa ci si aspettava, cosa ha funzionato, cosa no, quali problemi sono
emersi durante lo sviluppo e quali correzioni sono state applicate. L'obiettivo
\`e una discussione onesta, con ogni affermazione ancorata all'evidenza
sperimentale (Sez.~\ref{sec:results}).

\subsection{Obiettivi iniziali di ciascun esperimento}
\begin{itemize}[nosep]
  \item \textbf{CartPole controllo}: verificare che i metodi tabellari
        apprendano davvero e confrontare on-policy (SARSA) vs off-policy
        (Q\hbox{-}Learning).
  \item \textbf{Sensibilit\`a} ($\varepsilon,\alpha,\gamma$): mostrare l'effetto
        degli iperparametri sulla velocit\`a e stabilit\`a dell'apprendimento.
  \item \textbf{MC vs TD(0)}: illustrare il compromesso bias/varianza tra i due
        stimatori di predizione.
  \item \textbf{Casi di fallimento}: rendere visibili instabilit\`a
        ($\alpha$ estremo) e mancata esplorazione ($\varepsilon=0$).
  \item \textbf{Acrobot}: applicare gli stessi metodi a un ambiente pi\`u
        difficile (reward sparso, stato 6-D) e discutere i limiti della
        discretizzazione.
  \item \textbf{Pong/DQN}: implementare un DQN da zero e isolare, per ablazione,
        il contributo di experience replay e target network; testare Double DQN.
\end{itemize}

\subsection{Cosa ha funzionato}
\begin{itemize}[nosep]
  \item \textbf{Apprendimento tabellare corretto}: su CartPole entrambi gli
        algoritmi salgono da $\sim$30 fino a $\sim$175--225 (ancora in salita a
        5000 episodi dopo il fix \#2); Q\hbox{-}Learning mostra varianza tra seed
        pi\`u bassa (Tab.~\ref{tab:cartpole-control}).
  \item \textbf{Sensibilit\`a coerente con la teoria}: $\alpha=0.1$ \`e il
        compromesso migliore, $\gamma$ alto vince su un compito di
        sopravvivenza (Sez.~\ref{sec:res-cartpole-sens}).
  \item \textbf{Casi di fallimento riproducibili}: $\alpha=1.0$ e
        $\varepsilon=0$ producono i fallimenti attesi
        (Sez.~\ref{sec:res-cartpole-fail}).
  \item \textbf{Ablazione DQN da manuale}: replay indispensabile, target network
        acceleratrice; risultato netto e riproducibile
        (Fig.~\ref{fig:pong}, Tab.~\ref{tab:pong-ablation}).
  \item \textbf{Riproducibilit\`a}: media su 3 seed con bande, seeding combinato
        (numpy+torch+env), array grezzi salvati in \code{.npz} e tabelle
        generate da quei dati.
  \item \textbf{Bonus Double DQN}: implementato e funzionante (apprende in modo
        simile al DQN completo).
\end{itemize}

\subsection{Cosa non ha funzionato (o ha sorpreso)}
\begin{itemize}[nosep]
  \item \textbf{Il grafico MC vs TD(0) iniziale era poco informativo}: mostrava
        curve identiche, perch\'e tracciava il return della policy condivisa e
        non le stime $V(s)$. Corretto (fix \#1, Sez.~\ref{sec:lessons-fixes}).
  \item \textbf{Due ipotesi su Acrobot smentite dai dati}: bin pi\`u fini sono
        risultati \emph{migliori} (non peggiori) e $\varepsilon$ basso batte
        $\varepsilon$ alto (Sez.~\ref{sec:res-acrobot-bins},
        \ref{sec:res-acrobot-eps}). Risultati mantenuti e analizzati, non
        ``aggiustati''.
  \item \textbf{CartPole non raggiungeva 500} ma si fermava a $\sim$184: un
        \emph{tetto di discretizzazione}, non un bug. Il fix \#2 lo ha alzato
        ($\sim$225) ma al prezzo di convergenza pi\`u lenta e varianza maggiore
        --- un miglioramento \emph{a doppio taglio}, non netto (vedi
        Sez.~\ref{sec:lessons-fixes}).
  \item \textbf{Acrobot lontano dall'ottimo} ($\sim$$-345$ vs $\sim$$-100$) e con
        curve molto rumorose: atteso per reward sparso e spazio grande.
  \item \textbf{Double DQN non supera il DQN base} nel regime breve (200k step):
        la riduzione del maximization bias non \`e determinante con cos\`i pochi
        step; comportamento noto in letteratura, non un errore.
\end{itemize}

\subsection{Problemi incontrati durante lo sviluppo}
\label{sec:lessons-problems}
Questi sono problemi \emph{implementativi} reali (distinti dai limiti
intrinseci dell'RL discussi sotto):
\begin{itemize}[nosep]
  \item \textbf{GPU Blackwell (RTX 5060)}: richiede la build PyTorch CUDA 12.8;
        con build pi\`u vecchie la GPU non viene rilevata.
  \item \textbf{Pipeline Atari, tre trappole risolte}: (i) l'ambiente
        \code{-v5} ha gi\`a un frame-skip interno, che combinato con
        \code{AtariPreprocessing} darebbe uno skip di 16 frame $\Rightarrow$ si
        forza \code{frameskip=1} alla creazione; (ii) in Gymnasium 1.x
        \code{FrameStack} \`e stato rinominato \code{FrameStackObservation};
        (iii) i frame impilati arrivano gi\`a in formato channel-first
        \code{(4,84,84)}, quindi \emph{non} va applicata alcuna trasposizione
        prima della CNN. Un errore su uno qualsiasi di questi tre punti produce
        crash o input malformati alla rete.
  \item \textbf{Registrazione ALE}: con \code{ale-py} recente serve
        \code{gym.register\_envs(ale\_py)}, altrimenti \code{gym.make("ALE/...")}
        fallisce.
  \item \textbf{Memoria del replay buffer}: 1M transizioni Atari a frame interi
        sarebbero $\sim$28\,GB; si \`e limitato il buffer a 100k transizioni e si
        conservano i frame in \code{uint8}, normalizzando solo al training.
  \item \textbf{Backend Matplotlib}: forzato ad \code{Agg} (non interattivo) per
        salvare le figure senza GUI.
  \item \textbf{Conflitti di versione nel codice}: durante la gestione con
        \code{git} sono rimasti dei marcatori di merge non risolti in alcuni file
        (Atari/DQN/plotting). Sono stati risolti mantenendo la versione
        \emph{effettivamente eseguita} (quella che ha prodotto i risultati Pong),
        e verificati con un test di import dell'intero pacchetto.
\end{itemize}

\subsection{Correzioni applicate}
\label{sec:lessons-fixes}
\paragraph{Fix \#1 --- confronto MC vs TD(0) reso significativo.}
\emph{Cosa}: oltre al return della policy, si traccia ora la stima
$V(s_{\text{ref}})$ di uno stato di riferimento episodio per episodio.
\emph{Perch\'e}: il return \`e identico per costruzione tra i due metodi (stessa
policy), quindi non pu\`o mostrare la differenza; la differenza bias/varianza
vive nelle stime $V$.
\emph{Effetto}: su CartPole la stima MC salta rapidamente al valore corretto
($\sim$90 sulla policy appresa) con alta varianza iniziale, mentre TD(0) parte
da 0 e sale lentamente (distorsione da bootstrap) in modo liscio --- il quadro
classico di Sutton \& Barto \S6.2 (Fig.~\ref{fig:pred-value}).
\emph{Un bug scoperto e corretto durante questa verifica}: la prima scelta dello
stato di riferimento era il ``bin centrale'' su ogni dimensione. Su Acrobot ci\`o
d\`a uno stato \emph{fisicamente impossibile} --- l'osservazione codifica gli
angoli come $[\cos\theta, \sin\theta,\dots]$ e $\cos\theta,\sin\theta$ non
possono essere entrambi $\approx 0$ --- quindi quello stato non viene mai
visitato e la stima restava piatta a $0$. Corretto usando lo \emph{stato di
reset discretizzato}, sempre raggiungibile e visitato ad ogni episodio
(Fig.~\ref{fig:ac-pred-value}).
\emph{Risolto?} S\`i, entro l'approccio scelto (un singolo stato di
riferimento); una misura globale (RMSE di $V$ su tutti gli stati) richiederebbe
un $V$ di riferimento non disponibile qui.

\paragraph{Fix \#2 --- discretizzazione CartPole pi\`u efficiente.}
\emph{Cosa}: gli intervalli di clipping di posizione e angolo sono stati
ristretti dai limiti dello spazio di osservazione ($\pm4.8$, $\pm0.418$) ai
limiti di \emph{terminazione} ($\pm2.4$, $\pm0.21$).
\emph{Perch\'e}: oltre i limiti di terminazione l'episodio \`e gi\`a finito,
quindi met\`a dei bin di posizione e angolo era sprecata su stati che non
possono verificarsi.
\emph{Effetto (a doppio taglio, non netto)}: raddoppia la risoluzione utile sulle
due dimensioni pi\`u importanti e \emph{alza il tetto} --- le curve superano il
vecchio plateau di $\sim$184 e SARSA arriva a $\sim$225. Ma pi\`u stati
significano convergenza \emph{pi\`u lenta} (curve ancora in salita a 5000
episodi) e \emph{varianza maggiore} tra seed ($\pm 39$/$\pm 61$ contro
$\pm 8$/$\pm 18$). Inoltre l'ordine cambia: SARSA ora supera Q\hbox{-}Learning
(Tab.~\ref{tab:cartpole-control}).
\emph{Risolto?} Il ``tetto'' si alza ma non sparisce: resta un limite di ogni
metodo tabellare su stato continuo (serve approssimazione di funzione).

\subsection{Limiti residui}
\label{sec:lessons-limits}
Alcuni limiti restano anche dopo le correzioni. \`E importante distinguere quelli
\emph{intrinseci all'RL} da eventuali difetti di implementazione: i seguenti sono
caratteristiche attese, non bug.
\begin{itemize}[nosep]
  \item \textbf{Efficienza campionaria (sample efficiency) bassa}: i metodi
        tabellari richiedono migliaia di episodi (5k CartPole, 10k Acrobot) per
        risultati modesti. \`E inerente al tabular senza generalizzazione.
  \item \textbf{Stabilit\`a di convergenza}: le curve non si appiattiscono mai del
        tutto. Con $\varepsilon$ e $\alpha$ \emph{costanti}, l'agente continua a
        esplorare e pesa le esperienze recenti in modo esponenziale: la
        fluttuazione asintotica \`e attesa, non un errore.
  \item \textbf{Sparsit\`a del reward (Acrobot)}: reward $-1$/passo fino
        all'obiettivo rende l'assegnazione del credito difficile; spiega la
        convergenza lenta e la distanza dall'ottimo.
  \item \textbf{Sensibilit\`a agli iperparametri}: i risultati dipendono
        marcatamente da $\alpha,\gamma,\varepsilon$ e dal numero di bin (Sez.~
        \ref{sec:res-cartpole-sens}, \ref{sec:res-acrobot-bins}). Tipico dell'RL.
  \item \textbf{Varianza da seed / casualit\`a}: le bande $\pm1\sigma$ sono a
        volte ampie (Acrobot, predizione). Con solo 3 seed la stima
        dell'incertezza \`e essa stessa rumorosa; \`e un compromesso di budget,
        dichiarato apertamente.
  \item \textbf{Vincoli computazionali (Pong)}: 200k step contro i 50M del paper
        Nature; per questo Pong si ferma a $\sim$$-12$ e non a punteggi positivi.
        \`E una scelta di budget, non un difetto dell'implementazione.
  \item \textbf{Strategia di esplorazione}: $\varepsilon$-greedy costante \`e la
        pi\`u semplice; su Acrobot un valore alto \emph{costante} penalizza
        (Sez.~\ref{sec:res-acrobot-eps}). Un annealing sarebbe pi\`u adatto.
  \item \textbf{Generalizzazione / discretizzazione}: la discretizzazione
        uniforme perde precisione (tetto di CartPole; Acrobot lontano
        dall'ottimo). Non \`e overfitting in senso classico ma
        \emph{sotto-rappresentazione} dello stato: il limite \`e la griglia, non
        i dati.
  \item \textbf{Predizione, misura locale}: il fix \#1 mostra la varianza su
        \emph{un} stato di riferimento, non l'errore globale su $V$.
\end{itemize}

\subsection{Possibili miglioramenti futuri}
\label{sec:lessons-future}
Realistici, oltre gli obiettivi del corso:
\begin{itemize}[nosep]
  \item \textbf{Annealing di $\varepsilon$ e decadimento di $\alpha$}: curve pi\`u
        pulite e migliore convergenza asintotica.
  \item \textbf{Approssimazione di funzione per Acrobot}: tile coding o una
        piccola rete al posto della tabella, per superare la maledizione della
        dimensionalit\`a.
  \item \textbf{DQN pi\`u lungo e varianti}: pi\`u step, pi\`u Dueling DQN e
        Prioritized Experience Replay, verso punteggi positivi su Pong.
  \item \textbf{Predizione pi\`u rigorosa}: RMSE di $V$ rispetto a un riferimento
        Monte Carlo ad alta campionatura, come nell'esercizio \code{random\_walk}.
  \item \textbf{Ablazione Pong in funzione degli step ambiente}: confronto a
        parit\`a esatta di calcolo.
  \item \textbf{Fase di valutazione puramente greedy per il DQN}: separare il
        return di training (con esplorazione) da quello della policy finale.
\end{itemize}
